\begin{solapartat}
La reacció que ens demanen és:
\[ {}^{40}_{19}\,\mathrm{K} \rightarrow {}^{40}_{18}\,\mathrm{Ar} + {}^{0}_{1}\,\beta^{+} \quad \text{\punts{0,6}} \]
Com podem comprovar la partícula emergent és un positró \punts{0,4} (no és precís que comentin res respecte la possible producció de neutrins)
\end{solapartat}

\begin{solapartat}
La llei de desintegració la podem escriure com:
\[ N(t) = N_0\,\mathrm{e}^{-\frac{t\,\ln 2}{t_{1/2}}} \rightarrow m(t) = m_0\,\mathrm{e}^{-\frac{t\,\ln 2}{t_{1/2}}} \quad \text{\punts{0,2}} \]
\[ m(t = 5\cdot 10^{9}) = 10\ \mathrm{g}\ \mathrm{e}^{-\frac{5\cdot 10^{9}\,\ln 2}{1,25\cdot 10^{9}}} = 6,25\cdot 10^{-1}\ \mathrm{g} \quad \text{\punts{0,3}} \]
Tal com diu el enunciat, per saber l'edat de la pedra, hem de partir de la hipòtesi que tot el Ar de la roca prove de la desintegració del K, per tant en l'instant inicial teníem 20~g de K. \punts{0,1}, després de passar el temps $t$, en tenim 10~g, per tant hem de resoldre l'equació:
\[ 10\ \mathrm{g} = 20\ \mathrm{g}\ \mathrm{e}^{-\frac{t\,\ln 2}{t_{1/2}}} \quad \text{\punts{0,2}} \Rightarrow \]
\[ \frac{1}{2} = 2^{-\frac{t}{t_{1/2}}} \Rightarrow t = t_{1/2} = 1,25\cdot 10^{9}\ \text{anys} \quad \text{\punts{0,2}} \]
\end{solapartat}
